\documentclass[11pt]{article}

\usepackage[review]{acl}

\usepackage{times}
\usepackage{latexsym}
\usepackage[T1]{fontenc}
\usepackage[utf8]{inputenc}
\usepackage{microtype}
\usepackage{inconsolata}
\usepackage{graphicx}
\usepackage{booktabs}

\title{Non-Exchangeability in KV-Cache Relay: A Causal Audit of Multi-Agent LLM Pipelines}

\author{Anonymous ACL submission}

\begin{document}
\maketitle

\begin{abstract}
Efficient key-value (KV) cache relay between agents in multi-agent LLM
pipelines rests on an assumption that is rarely tested directly: that some
dimension of the cache -- which transformer layer, which channel, which
question's content -- can be treated as interchangeable for the purpose of
compression or reuse. We test this assumption across three structurally
distinct axes within a sequential three-agent (Reasoner--Verifier--Finalizer)
pipeline on multi-hop question answering, and find it false along all
three. Layers dropped for compression are not safely approximated by their
neighbors. Key vectors' channels cannot be freely mixed once rotary
position embeddings impose position-dependent structure on them: a
rotation-based quantization scheme that assumes otherwise collapses into
out-of-distribution output, while a channel-respecting fix recovers
accuracy statistically indistinguishable from our strongest layer-selection
baseline, at comparable compression and no calibration overhead. Two
results sit at the center of this paper. First, a causal audit --
substituting the relayed cache with zeroed, random, or another question's
real content -- shows a receiving agent's accuracy depends on the specific
transmitted content, not merely on receiving a non-empty cache: a
three-tier ordering (real content $>$ wrong-but-real content $>$ no real
content) that is statistically significant in every pairwise comparison and
holds under compression, not only in an idealized uncompressed setting.
Second, and more surprising: a calibration procedure's own internal
confidence does \emph{not} predict which of two model architectures is
safe to compress this way -- and points in the wrong direction in our
tests, with the more fragile architecture producing the more
confident-looking calibration signal. The failure that matters is
invisible from the inside. We term this pattern \textbf{non-exchangeability}
and position our contribution as a causally-grounded, cross-axis synthesis
of several 2026 findings on compression-specificity and
calibration-objective-dependence, reporting concretely where four
published KV-relay efficiency techniques transfer to a new deployment
setting and where they do not.
\end{abstract}

\section{Introduction}

Efficient key-value cache relay between agents in a multi-agent LLM
pipeline rests on an assumption that is rarely stated and, to our
knowledge, never directly tested: that some dimension of the cache -- which
transformer layer, which channel within a head, which question's content --
can be treated as generic or interchangeable for the purpose of
compression, reconstruction, or reuse. This paper tests that assumption
directly, across three structurally distinct axes, within one multi-agent
pipeline, and finds it false along all three: dropped transformer layers
are not safely approximated by their neighbors; key vectors' channels
cannot be freely mixed once rotary position embeddings have imposed
position-dependent structure on them; and a receiving agent's accuracy
depends on the specific content of the relayed cache, not merely on
receiving \emph{some} non-empty cache, as we show with direct causal
evidence.

We call this \textbf{non-exchangeability}: efficiency techniques for
KV-cache relay fail in proportion to how much they implicitly assume
exchangeability along an axis the underlying representation does not
actually support. The paper's sharpest single result is that this failure
is invisible from the inside: a calibration procedure's own confidence
does not predict which of two model architectures is safe to compress
along the depth axis, and fails in the \emph{wrong} direction -- the
architecture that collapses harder under layer-selection produces the
\emph{more} confident-looking signal (Section~\ref{sec:results}, Finding
6). We demonstrate the broader pattern by diagnosing and, where possible,
correcting four published KV-relay efficiency techniques within our
pipeline: a cross-agent offset-correction method \citep{kvcomm2025} fails
because it assumes a delta observed for one question transfers to
another; a rotation-based quantization scheme
\citep{turboquant2026,polarquant2026} fails specifically on key vectors
because it assumes head dimensions are interchangeable, disrupting RoPE's
position-dependent structure; a per-channel quantization fix
\citep{kivi2024} resolves this by respecting that structure instead; and
layer-selection's cost is architecture-dependent in a way its own
calibration signal does not predict. The content-identity finding is not
an artifact of testing only uncompressed relay: the identical three-tier
causal ordering replicates on both a layer-selected/quantized
configuration and a uniformly quantized one, more decisively than on
uncompressed relay -- the claim holds under the conditions a practitioner
would actually deploy.

Three recent papers anticipate pieces of this argument from different
angles -- that compression misses task-specific information
\citep{pitfallskv2026}, that layer redundancy depends on the calibration
objective rather than being a fixed property \citep{rethinkredundancy2026},
and that adaptive, non-fixed-ranking methods can generalize across
architectures where fixed ones might not \citep{adak2026}. We position
this paper's contribution as the causal, cross-axis synthesis these three
approach separately: direct causal validation (not inference from
compression metrics) on the content-identity axis, replicated evidence
across depth and position axes within one system, and a
falsified-predictor result showing the failure is invisible from a
technique's own confidence signal.

\section{Related Work}

\paragraph{Cross-agent KV-cache relay.} Relaying key-value cache instead
of decoded text between agents has been proposed to avoid redundant
prefill computation. KVCOMM \citep{kvcomm2025} introduced an anchor-based
framework that estimates and corrects KV-cache ``offset drift'' when
reusing cache across different prefix contexts, reporting over 70\% cache
reuse and up to 7.8$\times$ prefill speedup on retrieval-augmented
generation, math reasoning, and coding tasks. LatentMAS
\citep{latentmas2026} extends this to full latent-space collaboration,
exchanging autoregressively-generated latent ``thoughts'' via KV-cache
working memory rather than relaying already-decoded text, reporting token
savings and accuracy gains across nine math/science/code benchmarks. We
differ from both in relaying real, decoded reasoning text through a fixed
sequential three-agent chain (Reasoner $\rightarrow$ Verifier
$\rightarrow$ Finalizer) rather than offset-corrected reuse or latent
thought generation, and in evaluating on multi-hop QA
\citep{hotpotqa2018}, which neither tested on.

\paragraph{Compressing the relayed cache.} Reducing what must be
transmitted between agents is a natural complement to relay itself. We
adopt two families of technique: layer-wise selection, which drops
calibration-identified low-importance transformer layers before
transmission and reconstructs them at the receiver via \texttt{zeros},
\texttt{nearest}, or \texttt{interpolate}, and per-head quantization,
which reduces numeric precision. A third reconstruction strategy,
Orthogonal Backfill \citep{orthobackfill2026}, injects a low-rank residual
of discarded content orthogonal to what is retained rather than
substituting or discarding it outright; we considered but did not
implement it (see Limitations for why). KV-cache compression is also
under active industrial development: KVTC \citep{kvtc2026} applies
media-compression-style transform coding -- PCA decorrelation, adaptive
quantization, entropy coding -- for up to 20$\times$ compression, a third
paradigm distinct from uniform quantization (ours, KIVI) and rotation
(TurboQuant/PolarQuant); NVIDIA's kvpress library and TensorRT-LLM's
reuse-aware cache confirm this is an actively contested axis of real
deployed systems, not a narrow academic concern.

\paragraph{Rotation-based quantization interacts badly with RoPE-encoded
keys -- a finding, not just an implementation note.} We initially
implemented a rotation-based quantization scheme in the style of
TurboQuant \citep{turboquant2026} and PolarQuant \citep{polarquant2026} --
two distinct papers sharing only two authors, not one work under two
names -- applying a fixed orthonormal Hadamard rotation to both keys and
values before quantizing to redistribute per-head outlier magnitude
(covering PolarQuant's rotate-then-quantize step, not TurboQuant's
separate QJL bias-correction term). On our real model this produced
severely out-of-distribution decoded output, not the graceful degradation
a generic rotation-based scheme predicts. We traced this to keys' rotary
position embeddings (RoPE): K is cached \emph{after} RoPE is applied,
which mixes channel pairs by a position-dependent phase, and a second,
RoPE-agnostic rotation on top disrupts that structure rather than merely
redistributing outlier magnitude. An isolating ablation (rotating only V,
leaving K unrotated) reverted the failure to ordinary quantization-noise
degradation, confirming the interaction was with K specifically --
consistent with an established interaction in the KV-quantization
literature (next paragraph), and reported here as a diagnosed negative
result, not an implementation detail.

\paragraph{RoPE-aware KV cache quantization.} A separate line of work
addresses exactly this interaction directly. KIVI \citep{kivi2024}
quantizes the key cache per-channel and the value cache per-token --
motivated by RoPE leaving channel-wise magnitude more consistent across
positions than uniform per-head ranges assume -- rather than by rotating.
RotateKV \citep{rotatekv2025} instead applies rotation \emph{before} RoPE
via Pre-RoPE Grouped-Head Rotation, explicitly designed to avoid the
interaction we observed. KVQuant \citep{kvquant2024} and related work
similarly identify post-RoPE key quantization as harder than value
quantization for exactly this reason. We adopt a scoped, KIVI-inspired fix
(per-channel key quantization only, no group-wise windowing or streaming
residual buffer from the original method) and find it resolves the
collapse -- 0.0\%/0.0\% to 52.0\%/64.1\% F1 at $n=100$, statistically
indistinguishable from our strongest layer-selection config (\texttt{D};
McNemar $p=0.86$) at \texttt{B\_int8}-level compression ($\approx$2.00x)
with no calibration profile required -- direct empirical confirmation, on
our pipeline and task, of what this literature identifies as the correct
axis for addressing the problem, rather than the rotation-based axis we
tried first. Whether KIVI's other asymmetric half (per-token value
quantization) adds anything on top is tested in Results, Finding 4.

\paragraph{Auditing whether KV reuse does what it claims.} A recent line
of work interrogates cross-agent KV/latent reuse critically rather than
only reporting end-task accuracy. ``When KV Cache Reuse Fails in
Multi-Agent Systems'' \citep{kvreusefails2026} shows reuse strategies
effective for generation agents can silently corrupt an LLM judge's
cross-candidate comparison even when end-task accuracy appears stable.
``When Does Latent Communication Pay?'' \citep{latentcommpay2026} proposes
replacing relayed KV with mismatched, zeroed, or moment-matched-random
substitutes to test whether accuracy is causally attributable to the
specific transmitted content rather than to the receiver simply having a
non-empty cache -- a methodology we adopt directly for our own causal
audit (Section~\ref{sec:method}). We extend this tradition along an axis
neither paper covers: cross-architecture generalization, treating ``does
the same technique, faithfully reproduced, transfer to a different model
family'' as itself a target of audit.

\paragraph{Layer redundancy is not universal.} ``No Free Swap''
\citep{nofreeswap2026} finds that which transformer layers are
functionally interchangeable across Qwen3-8B, Llama-3.1-8B, and
Mistral-7B-v0.1 depends on the model and evaluation protocol rather than
being a fixed architectural property. This is consistent with our own
finding (Section~\ref{sec:results}, Finding 3) that a near-identical
proportion of dropped layers costs substantially more accuracy on
Mistral-7B-Instruct-v0.3 than on Qwen2.5-7B-Instruct, and offers one
candidate mechanism for why.

\paragraph{Three 2026 papers anticipate pieces of this paper's argument.}
``Rethinking Layer Redundancy'' \citep{rethinkredundancy2026} argues
redundancy is a joint function of model \emph{and calibration objective},
and that ``a universal layer ranking may not exist'' -- raising the
question of whether our cross-architecture gap is really about the model,
or about our calibration ranking failing to transfer. We tested the most
direct proxy: whether the calibration signal's own confidence (the
separation between kept- and dropped-layer importance scores) predicts
which model is safe to compress. It does not, and fails in the
\emph{wrong} direction -- Mistral's tier separation (0.673) is larger, not
smaller, than Qwen's (0.491), despite Mistral being the more fragile model
when that ranking is acted upon. This is a direct empirical test of
exactly the concern that paper raises, and supports a stronger reading:
not just that redundancy is calibration-objective-dependent, but that the
objective's own confidence is not a reliable signal of when its ranking is
safe to trust -- consistent with our broader claim that non-exchangeability
is not visible from a technique's internal signals.

AdaK \citep{adak2026} reports that adaptive KV-budget estimation
generalizes across Qwen3 (4B/8B) and Mistral-7B-Instruct-v0.2 -- the same
two architecture families this paper tests, not a third (Llama is not
among AdaK's evaluated models, though a superficial reading of the
abstract could suggest otherwise) -- which on its face looks like a
counterexample to our architecture-dependence claim. We do not believe it
is: AdaK estimates budget per-instance and per-model at inference time
rather than committing to a single calibration-time ranking applied
uniformly thereafter, so it never assumes a fixed, transferable ranking is
safe to act on. Read this way, AdaK's success is a positive instance of
our principle, not a counterexample -- it works specifically because it
avoids the fixed-ranking assumption our own \texttt{D} configuration
makes. This is our reading of the mechanism, not something we verified by
reimplementing AdaK ourselves; a direct head-to-head comparison on our own
pipeline is noted as a natural extension in Limitations.

Finally, ``The Pitfalls of KV Cache Compression'' \citep{pitfallskv2026}
makes a general version of the point this paper investigates concretely:
that compression methods optimize for aggregate metrics while implicitly
equating token retention with functional preservation. Our contribution is
not the abstract observation but a causally-grounded, cross-axis
demonstration of it within one system -- the same failure pattern
recurring across three structurally distinct axes of KV-cache structure
(depth, position, content-identity) in a single pipeline, with the
content-identity axis causally validated via our audit rather than
inferred from compression metrics alone. We position this paper as the
causal, cross-axis synthesis these three papers each gesture toward from a
different angle.

\section{Method}
\label{sec:method}

\paragraph{Pipeline.} We study a sequential three-agent pipeline --
Reasoner, Verifier, Finalizer -- answering HotpotQA
\citep{hotpotqa2018} (distractor configuration) questions, where each item
bundles its own ten candidate passages and requires no retrieval step;
Section~\ref{sec:results}'s GSM8K \citep{gsm8k2021} generalization check
uses the same pipeline structure. Agents share the same prompt structure
and role definitions across every configuration in this paper; the only
thing that varies between configurations is \emph{how} one agent's output
reaches the next. In \texttt{text\_agent}, the Reasoner's decoded text is
passed to the Verifier as a literal string, which reads and re-processes
it from scratch. In every KV-relay configuration (\texttt{A} through
\texttt{E}), the Reasoner's key-value cache -- its internal representation
of the question, context, and its own generated reasoning -- is passed
directly to the Verifier, which continues generation from that cache
instead of re-reading text. \texttt{single\_agent} is a no-relay control:
one model call, no intermediate agents at all. All KV-relay configurations
share identical prompts with \texttt{text\_agent}, so any accuracy or
latency delta between a KV configuration and \texttt{text\_agent} is
attributable to the relay mechanism itself, not to prompt differences.

\paragraph{Models.} We evaluate on Qwen2.5-7B-Instruct
\citep{qwen2p5_2024} (bf16) as our primary model, and -- specifically to
test whether findings generalize across architecture, not just within one
model family -- Mistral-7B-Instruct-v0.3 \citep{mistral7b2023} (bf16),
chosen for its architectural differences (distinct layer count, attention
head configuration, and tokenizer) at comparable parameter scale. Both are
decoder-only transformers using grouped-query attention and RoPE.

\paragraph{Compression configurations.} \texttt{B\_int8}/\texttt{B\_int4}
apply uniform per-head min-max quantization to every transmitted layer's
key and value tensors (a fix to \texttt{B\_int4}'s originally-catastrophic
failure is detailed below). \texttt{C} applies calibration-driven layer
selection -- dropping transformer layers a per-model calibration pass
identifies as low-importance before transmission, reconstructing dropped
layers at the receiver via one of three strategies (\texttt{zeros},
\texttt{nearest}, \texttt{interpolate}). \texttt{D} combines layer
selection with adaptive per-layer-tier quantization (high-importance
layers at 8-bit, medium-importance at 4-bit). \texttt{E} adds an
anchor-based cross-agent offset-correction mechanism adapted from KVCOMM
(Finding 2 is a negative result for this configuration).

\paragraph{RoPE-aware quantization fix for \texttt{B\_int4}.}
\texttt{B\_int4} (uniform 4-bit, no layer selection) initially collapsed
to 0.0\% accuracy -- diagnosed as a quantization/RoPE interaction and
confirmed via two ablations (see Related Work for the full diagnostic
story). Our fix (\texttt{B\_int4\_kivi}), inspired by KIVI, quantizes keys
per-channel (one numeric range per head dimension, computed across
positions) instead of per-head, leaving values unchanged.

\paragraph{Statistical methodology.} For every headline comparison between
two configurations, we report exact (binomial-form) McNemar's test on
paired per-example correctness -- appropriate because both configurations
are evaluated on the identical example set -- and a paired bootstrap 95\%
confidence interval (10,000 resamples) for the difference in mean
token-F1. We additionally report Monte Carlo-estimated statistical power
at the current sample size for comparisons that do not reach significance,
so that ``not significant'' and ``confirmed no effect'' are never
conflated in our reporting.

\paragraph{Causal audit.} To test whether a configuration's accuracy
reflects the receiving agent using the specific transmitted content,
rather than merely having \emph{some} non-empty cache to attend over, we
substitute the relayed KV cache with (i) an all-zero cache, (ii) a
moment-matched random cache (same per-tensor mean/standard deviation, no
real structure), and (iii) another, unrelated held-out question's real
cache at the same pipeline stage (sampled from a pool built from examples
never included in the scored evaluation set, so no scored question's
content can leak back to itself). We apply this to our three strongest
configurations (\texttt{A}, \texttt{B\_int8}, \texttt{D}).

We deliberately do not treat these three conditions as interchangeable
evidence for the same claim. Zero-ablation is a known-imperfect causal
baseline in the interpretability literature: it pushes activations off the
training distribution, which can produce large behavioral changes for
reasons unrelated to whether the zeroed content carried meaningful
information. We treat the moment-matched-random and mismatched-example
conditions -- which preserve realistic activation statistics or substitute
another real, in-distribution cache -- as the more informative evidence,
and zero-ablation as a supplementary sanity floor; if zeroing alone
diverges sharply from the other two, we attribute that to the
out-of-distribution effect rather than content-dependence. A fully
rigorous treatment of ``no difference'' would also use a pre-registered
equivalence margin rather than absence of significance under our paired
test; we did not have the budget to pre-register and power such a test,
and report significance/CI results with that caveat.

\section{Results}
\label{sec:results}

Eight findings follow. Two are the paper's headline results: Finding 5,
the causal audit establishing that relayed KV carries specific, real
content rather than merely being non-empty, and Finding 6, that a
calibration procedure's own confidence fails to predict -- and actively
misdirects on -- which architecture is safe to compress. The remaining
findings characterize specific published techniques against this
backdrop.

\begin{table*}[t]
\centering
\small
\begin{tabular}{lcccc}
\toprule
\textbf{Config} & \textbf{Accuracy} & \textbf{F1} & \textbf{Latency} & \textbf{KV/hop} \\
\midrule
single\_agent & 57.0\% & 68.4\% & 1.6s & --- \\
text\_agent & 54.0\% & 68.3\% & 7.2s & --- \\
A & 56.0\% & 69.1\% & 8.6s & 144.35 MB \\
B\_int8 & 57.0\% & 70.5\% & 8.3s & 72.16 MB (2.00x) \\
B\_int4 (original) & 0.0\% & 0.0\% & 39.4s & $\sim$76.8 MB (2.00x) \\
\textbf{B\_int4\_kivi (fixed)} & \textbf{52.0\%} & \textbf{64.1\%} & 9.6s & $\sim$72.6 MB (2.00x) \\
C & 43.0\% & 57.8\% & 10.0s & 104.16 MB \\
D & 50.0\% & 61.9\% & 10.2s & $\sim$52.2 MB (2.77x vs A) \\
E & 9.0\% & 17.4\% & 19.8s & $\sim$49.1 MB (2.00x) \\
E\_int8 & 1.0\% & 6.2\% & 22.2s & 49.05 MB \\
\bottomrule
\end{tabular}
\caption{Established results (Qwen2.5-7B-Instruct, HotpotQA, $n=100$).
Ratios are compressed-vs-A's 144.35~MB, except D (which drops layers before
compressing) shows its ratio vs.\ A rather than its own self-referential
2.00x, since vs.-A is the number a reader wants. See Limitations for why
4-bit and 8-bit configs share the same $\approx$2.00x ceiling here.}
\label{tab:main}
\end{table*}

\paragraph{Finding 1 -- uniform 8-bit KV quantization is statistically
indistinguishable from uncompressed relay.} \texttt{A} vs \texttt{B\_int8}:
McNemar $p=1.0$, only 3 discordant examples out of 100. Confirmed on a
second model family (Mistral-7B-Instruct-v0.3, $n=50$): 0 discordant
examples -- identical per-example outcomes. This is not an underpowered
null result; the observed effect is tight enough at this sample size to be
a genuine near-zero cost.

\paragraph{Finding 2 -- KVCOMM-style offset correction (\texttt{E}) fails
catastrophically, independent of a decoding confound we explicitly ruled
out.} \texttt{D} vs \texttt{E}: McNemar $p\approx0.00003$ (both before and
after an unrelated repetition-penalty fix that changed several other
configurations' numbers -- see Limitations). Manual inspection of
\texttt{E}'s raw output shows duplicated phrases and hallucinated content
inconsistent with the anchor-corrected cache carrying usable information,
not merely a lower-quality-but-coherent answer.

\paragraph{Finding 3 -- layer-selection's accuracy cost is
architecture-dependent.} On Mistral-7B-Instruct-v0.3 ($n=50$), dropping
9/32 layers (28\%) cost 30.5 F1 points relative to uncompressed relay
($p=0.0001$). On Qwen2.5-7B-Instruct ($n=100$), a near-identical
proportion (8/28, 29\%) shows a smaller, not-yet-significant cost
($p=0.34$) -- the cross-model claim does not rest on that reaching
significance, though: even Qwen's most generous uncertainty bound (F1 CI
upper +16.4 points) is under half of Mistral's confirmed collapse.
Consistent with prior work finding layer redundancy architecture- and
protocol-dependent, not a fixed property of a given compression ratio.

\paragraph{Finding 4 -- \texttt{B\_int4}'s fixed variant matches our
strongest layer-selection configuration's accuracy at
\texttt{B\_int8}-level compression and lower complexity -- but not at
higher compression than \texttt{D}.} \texttt{B\_int4\_kivi}
(52.0\%/64.1\%) vs \texttt{D} (50.0\%/61.9\%): McNemar $p=0.86$, F1 delta
+2.2 points [95\% CI: $-8.2$, $+12.8$] -- statistically indistinguishable
-- and requiring no per-model calibration profile. On compression: this
codebase's quantized tensors are stored as \texttt{torch.uint8} regardless
of bit-width (no nibble-packing for either 4-bit or 8-bit; see
Limitations), so \texttt{B\_int4\_kivi}'s real footprint is
$\approx$72.6~MB/hop ($\approx$2.00x vs \texttt{A}), matching
\texttt{B\_int8} -- not higher compression than \texttt{D}, whose real
footprint is $\approx$52.2~MB/hop ($\approx$2.77x vs \texttt{A}).
\textbf{\texttt{D} produces a $\sim$39\% smaller cache than
\texttt{B\_int4\_kivi}}, because \texttt{D}'s compression advantage is
layer-selection's savings from dropping layers entirely, not int4
precision. \texttt{B\_int4\_kivi}'s defensible claim is therefore: it
matches \texttt{D}'s accuracy at \texttt{B\_int8}-level compression with
no calibration profile -- a real, useful simplicity trade-off, not a
compression \emph{win} over \texttt{D}. It also trends below
\texttt{A}/\texttt{B\_int8} on accuracy (52\% vs 56--57\%), but this gap
is not yet statistically confirmed at $n=100$ (7--12 discordant examples,
$p=0.36$--$0.50$); we report this as an open question, not as cost-free
the way Finding~1 establishes for \texttt{B\_int8}. Also fixing the value
side (\texttt{B\_int4\_kivi\_full}, per-token quantization) does not
improve on this further (44.0\% vs 50.0\% on a matched 50-example subset,
$p=0.51$, trending the other way) and, even after the same correction, is
still worse on compression ($\approx$1.94x vs \texttt{B\_int4\_kivi}'s
$\approx$2.00x) for a real, mechanical reason unrelated to the
byte-accounting bug: per-token quantization stores a scale/zero-point pair
per sequence position, and sequence length (hundreds to low thousands of
tokens) vastly exceeds K's per-channel grouping (128 channels) -- real
overhead, no accuracy benefit. Together with an earlier null result for
rotating values instead, two independently-motivated value-side
interventions both failed -- evidence the key/RoPE interaction was the
entire mechanism behind \texttt{B\_int4}'s collapse, not one of several
factors.

\begin{figure}[t]
\centering
\includegraphics[width=\columnwidth]{figures/fig1_causal_audit.png}
\caption{Causal audit: real $>$ mismatched $>$ zeroed $\approx$ random,
replicated across \texttt{A}, \texttt{D}, and \texttt{B\_int8}. Every
pairwise comparison is statistically significant.}
\label{fig:causal_audit}
\end{figure}

\paragraph{Finding 5 -- headline result: relayed KV demonstrably carries
specific, real content, not merely a non-empty cache.} We ran the causal
audit (Section~\ref{sec:method}) on configuration \texttt{A} (uncompressed
relay), $n=50$: \texttt{A} reaches 50.0\% accuracy / 64.0\% F1;
\texttt{A\_audit\_zeroed} and \texttt{A\_audit\_random} both collapse to
0.0\%/0.0\%; \texttt{A\_audit\_mismatched} (a real, different held-out
question's cache) lands at 28.0\%/37.1\% -- between the two. Every
pairwise comparison in this three-tier ordering is statistically
significant: \texttt{A} vs zeroed and \texttt{A} vs random, McNemar
$p<0.0001$ each; \texttt{A} vs mismatched, $p=0.0127$; mismatched vs
zeroed and mismatched vs random, $p=0.0001$ each (an earlier $n=20$ pilot
showed the same ordering with the \texttt{A}-vs-mismatched comparison just
short of significance, $p=0.0625$, resolved at $n=50$). This is, to our
knowledge, the first direct causal validation of this kind for a
sequential (not fan-in) multi-agent KV-relay chain, and it underwrites
every compression result in this paper: Findings 1 and 4 (near-zero-cost
8-bit quantization, and the \texttt{B\_int4} fix matching \texttt{D}) are
only meaningful claims about \emph{preserved information} if the
underlying relay is shown to carry real content in the first place, which
this finding establishes directly rather than assuming.

Qualitative inspection adds a smaller, mechanistically interesting
observation: \texttt{A\_audit\_mismatched}'s incorrect answers are
coherent, grammatical, and plausible -- real linguistic structure attached
to the wrong question -- while \texttt{A\_audit\_zeroed} produces garbled
but recognizable English, and \texttt{A\_audit\_random} produces content
qualitatively \emph{more} degraded (code-fragment tokens, mixed-language
noise) despite matching the real cache's per-tensor statistics. We
speculate this reflects how attention treats the two null conditions
differently: a zero vector is a weak key attention can largely
down-weight, while realistic-magnitude random noise resembles genuine
signal and can actively misdirect attention rather than being ignored. A
secondary, exploratory observation, not a claim we isolated further.

\paragraph{Finding 5, extended -- the same ordering holds under
compression and on a second task family.} Extending the audit to
\texttt{D} (layer-selected + quantized) and \texttt{B\_int8} (uniformly
quantized), $n=50$ each on HotpotQA, the identical three-tier ordering
replicates more decisively than on uncompressed \texttt{A}: \texttt{D}
54.0\%/60.2\% vs. \texttt{D\_audit\_zeroed}/\texttt{D\_audit\_random}
$\approx$0.0\% vs. \texttt{D\_audit\_mismatched} 26.0\%/32.2\%;
\texttt{B\_int8} 54.0\%/68.0\% vs. zeroed/random 0.0\%/0.0\% vs.
mismatched 24.0\%/34.4\%. All ten pairwise comparisons across both
configurations are significant (real vs.\ zeroed/random $p<0.0001$ for
both; real vs.\ mismatched $p=0.0013$ \texttt{D} / $p=0.0003$
\texttt{B\_int8}, both tighter than \texttt{A}'s $p=0.0127$; mismatched
vs.\ zeroed/random $p=0.0002$--$0.0005$), and raw generations show the
identical failure signatures found on \texttt{A}. A single further
confirmatory run on GSM8K \citep{gsm8k2021} (Qwen, $n=50$, \texttt{A} and
the full causal audit on \texttt{A}) shows the same pattern still more
decisively: \texttt{A} 90.0\%, \texttt{A\_audit\_zeroed} 0.0\%,
\texttt{A\_audit\_random} 2.0\%, \texttt{A\_audit\_mismatched} 28.0\%,
every comparison significant ($p<0.0002$); the same run's \texttt{D}
matched \texttt{A}'s accuracy exactly (McNemar $p=1.0$) --
layer-selection was statistically free on this task, unlike on HotpotQA.
Together these ground the content-identity claim for compressed,
layer-selected relay and for a second task family, not only an idealized
uncompressed HotpotQA baseline -- though the GSM8K check is one run on one
model, reported as suggestive, not a systematic cross-task study (see
Limitations).

\paragraph{Finding 6 -- headline result: a calibration signal's own
confidence does not predict downstream layer-selection safety, and fails
in the wrong direction.} We tested the most direct available proxy for
whether our cross-architecture gap (Finding 3) reflects the model or
merely our calibration procedure: the separation between
calibration-importance scores assigned to kept (tier-1/2) versus dropped
(tier-3) layers, as a measure of how confidently the calibration ranks
layers. Qwen's mean separation is 0.491 (tier-1 0.737, tier-3 0.246);
Mistral's is \emph{larger}, 0.673 (tier-1 0.851, tier-3 0.178) --
Mistral's calibration looks more decisive, not less, yet Mistral is the
model whose accuracy collapses harder when that ranking is acted upon
(Finding 3). This rules out a noisier or less-confident calibration signal
for Mistral as the explanation for Finding 3, and supports a sharper one:
the calibration objective can be equally or more internally confident
while being less reliable, undetectable from the calibration output
alone. This is a single comparison across two models, not a validated
general predictor (see Limitations), but a real, falsified hypothesis
test, not an assumed conclusion.

\paragraph{Finding 7 -- architecture-dependent layer-selection failures
differ in kind, not only in rate.} Beyond the accuracy gap already
reported (Finding 3), \texttt{D}'s incorrect answers differ in character
between models. Qwen's are short (mean 18 characters, 2\% exceed 150) and
overwhelmingly close, traceable near-misses. Mistral's are markedly longer
(mean 109 characters, 15\% exceed 150, maximum 1,454) and include a
genuine long tail of severe failures -- fluent but entirely fabricated
tangents, degenerate repetition loops -- absent from Qwen's failure set.
We checked an obvious alternative explanation -- under-penalizing
repetition for Mistral relative to its own tuned defaults -- against each
model's published generation defaults, and find the opposite: Qwen's own
default \emph{is} our \texttt{repetition\_penalty=1.05}, while Mistral's
specifies none, so our setting applies \emph{more} correction for
Mistral, not less. The long-tail pattern persists anyway, weakening a
decoding confound as the explanation and supporting depth-axis
non-exchangeability differing in \emph{character}, not only magnitude,
across architectures.

\begin{figure}[t]
\centering
\includegraphics[width=\columnwidth]{figures/fig2_b_int4_journey.png}
\caption{The \texttt{B\_int4} diagnosis-to-fix journey: collapse, a broken
rotation-based attempt, an isolating diagnostic, the working fix
(\texttt{B\_int4\_kivi}), and two value-side variants that did not
improve on it.}
\label{fig:b_int4_journey}
\end{figure}

\paragraph{Finding 8 -- donor-question content occasionally, but rarely,
bleeds through under mismatched substitution.} Unlike ``When Latent
Agents Lie'' \citep{latentagentslie2026}, which studies an adversarial
agent deliberately substituting hidden state in a fan-in topology, our
substitution is a non-adversarial intervention in a sequential chain: when
the receiving agent is wrong, does its answer reflect the donor's
content, or ordinary confusion on the real question? Manually reviewing
20 of 36 incorrect \texttt{A\_audit\_mismatched} answers against their
donor question, we find one unambiguous instance -- a Kansas university
fight-song question produced ``Ellie Goulding'' in its wrong answer,
traceable only to a donor question about that singer -- and one weaker,
ambiguous case; the remainder show ordinary failure on the real question's
own topic. A genuine but rare phenomenon, not the dominant explanation for
\texttt{A\_audit\_mismatched}'s accuracy loss -- a manual review of a
subset, not an exhaustive classification (see Limitations).

\section*{Limitations}

\paragraph{Causal audit and bleed-through scope.} We ran the causal audit
on \texttt{A} (uncompressed relay), \texttt{D} (layer-selected +
quantized), and \texttt{B\_int8} (uniformly quantized) -- the three-tier
ordering replicates across all three, at $n=50$ each -- but not on
\texttt{C} or the \texttt{B\_int4} family (including the
\texttt{B\_int4\_kivi} fix); we claim the mechanism confirmed for the
three tested configurations, not for every configuration in this paper,
though extending it is a natural next step given the consistency already
observed across three architecturally distinct conditions (no
compression, compression only, compression plus layer selection). The
bleed-through analysis (Finding 8) is narrower still: a manual review of
20 of 36 incorrect \texttt{A\_audit\_mismatched} examples by one
annotator, with no automated classifier, inter-annotator agreement check,
or statistical test on the resulting rate. We report it as a qualitative,
exploratory observation -- evidence that donor bleed-through can occur,
illustrated by one unambiguous example -- not a validated estimate of how
often it occurs; a rigorous version would require an automated or
multiply-annotated classification over the full set of incorrect
examples, which we did not have the budget to complete. Topology is also
fixed: every result in this paper is on a sequential, non-adversarial
chain. \citet{latentagentslie2026} audits an adversarial fan-in topology
instead, and LatentMAS \citep{latentmas2026} collaborates via layer-wise
KV concatenation across a different multi-agent structure than our
hop-by-hop handoff; whether the same three-tier causal ordering holds
under a fan-in or concatenation-based topology is a natural, cited
extension we did not attempt, not one we found and characterized.

\paragraph{Single-process, single-GPU evaluation.} Every experiment runs
within one Python process on one RTX 4090: no \texttt{KVMessage} is ever
serialized or transmitted over a network. Compression-ratio and
transmitted-byte accounting reflects what a real distributed deployment
would need to send and is meaningful on that basis, but our latency
numbers exclude network transfer time, and smaller relayed caches are not
observed to be faster in this setup, since the model still executes every
transformer layer regardless of what its injected cache contains -- claims
here characterize communication-cost tradeoffs and accuracy effects, not
measured end-to-end deployment speedups. This single-GPU constraint also
bounded the sample sizes and the number of model/config combinations we
were able to test; a systems-scale deployment study (multi-GPU, real
network transfer, throughput under concurrent load) was out of scope.

\paragraph{Statistical power and replication.} We report paired
significance tests (exact McNemar) and bootstrap confidence intervals for
every headline comparison; several -- INT8's near-zero cost, offset
correction's catastrophic failure, Mistral-7B's layer-selection collapse
-- are significant at conventional thresholds even at our evaluation
scale. Others are not: at $n=100$, we cannot yet statistically distinguish
Qwen2.5-7B's layer-selection condition from its uncompressed baseline in
isolation (power $\approx$15\%; roughly 3$\times$ the current sample size
would be needed). Our cross-architecture claim rests on the independently
well-powered Mistral-7B result, not the underpowered Qwen-only
comparison. Relatedly, \texttt{B\_int4\_kivi}'s comparison against
\texttt{D} (statistically indistinguishable, McNemar $p=0.86$) rests on a
single $n=100$ run each, not a multi-seed or replicated evaluation -- we
report the confidence interval alongside the point estimate specifically
so ``matches \texttt{D}'' is read as the current best estimate under one
run, not a claim strengthened by replication we have not performed.

\paragraph{Scope: two models, mostly one task, greedy decoding.} We test
Qwen2.5-7B-Instruct and Mistral-7B-Instruct-v0.3 (chosen for
architectural diversity at comparable scale); whether findings --
particularly layer-selection's architecture-dependence -- extend to
substantially larger or smaller models, mixture-of-experts architectures,
or models outside this decoder-only, grouped-query-attention, RoPE
lineage is untested (additional model families for this extension are in
progress at time of writing; one methodological risk worth flagging in
advance: \citet{latentcommpay2026} reports at least one model pairing
where the receiver could not consume a relayed cache at all, i.e.
cross-architecture relay can fail at the precondition stage, not just on
accuracy -- a null result under those conditions would itself be
informative, not a bug). All of our own architecture variation is
\emph{between} pipeline runs (the same model family fills every agent
role within a given run); whether a single pipeline mixing architectures
across agent roles -- a Qwen Reasoner handing its cache to a Mistral
Verifier, for instance -- is even viable is untested and, to our
knowledge, an open problem in this literature more broadly: KVCOMM itself
flags relay between ``agents with identical architectures but different
weights'' and agents with ``different attention formulations'' as work it
leaves for future exploration, and we inherit that same gap rather than
closing it. Nearly all results are on HotpotQA
(distractor configuration); we ran one $n=50$ confirmatory check on GSM8K
(Qwen, \texttt{A}/\texttt{D}/full causal audit on \texttt{A}) and both
headline findings replicated cleanly -- \texttt{D} matched \texttt{A}'s
accuracy exactly (McNemar $p=1.0$, more decisively free than on
HotpotQA) and the causal audit's three-tier ordering held with every
pairwise comparison significant ($p<0.0002$), larger effect sizes than the
original HotpotQA run. This is one run on one additional task family, not
a systematic multi-task study -- we did not extend it to
\texttt{B\_int8}/\texttt{D}-family causal audits or Mistral, and report it
as suggestive, not as evidence the claims generalize broadly across
tasks. This matters specifically for the offset-correction negative
result: KVCOMM was evaluated on MMLU, GSM8K, and HumanEval, none of which
is multi-hop QA, and whether the technique fares differently on the more
homogeneous task family it was originally validated on remains open (our
own GSM8K check did not re-test \texttt{E}). All results additionally use
greedy decoding (\texttt{do\_sample=False}); we did not evaluate under
sampling or self-consistency, which could interact differently with a
degraded or compressed relayed context than with a full one.

\paragraph{No bit-packing: reported compression ratios reflect
quantization granularity, not physical storage savings beyond
bf16$\rightarrow$8-bit.} Every quantized tensor in this codebase -- 4-bit
and 8-bit alike -- is stored as \texttt{torch.uint8}; no nibble-packing is
implemented. A 4-bit value therefore occupies the same one byte an 8-bit
value does, so any config whose transmitted layers are entirely quantized
achieves the same $\approx$2.00x compression relative to bf16 regardless
of nominal bit-width (see Table~\ref{tab:main}'s caption and Finding 4); a
version with real nibble-packing would compress 4-bit layers roughly twice
as far as reported here. \texttt{D}'s real compression advantage over
\texttt{B\_int8}/\texttt{B\_int4\_kivi} is entirely attributable to layer
selection dropping layers outright, not to its adaptive quantization
tier; implementing genuine 4-bit packing is future work that would benefit
every quantized configuration in this paper equally, not change any of its
comparative claims.

\paragraph{Partial reproduction of adapted and considered techniques.}
Our rotation-based quantization covers PolarQuant's core
rotate-then-quantize mechanism but omits TurboQuant's QJL bias-correction
term; our KIVI-inspired fix (Section~\ref{sec:method}) omits the original
method's group-wise windowing and fp16 residual buffer, implementing
full-sequence per-channel quantization only -- the minimal change
isolating the RoPE-interaction mechanism we diagnosed, not a complete
reproduction. We report both as inspired-by implementations scoped to
test a specific hypothesis, not certified reproductions, and encourage
readers wanting the original techniques' full guarantees to consult the
primary sources. We considered but did not implement a fourth
layer-reconstruction strategy, Orthogonal Backfill
\citep{orthobackfill2026}: its published formula requires attention
weights this pipeline's \texttt{sdpa} decode path does not expose, and
given our own evidence that naive/simplified reproductions can actively
mislead (the rotation-based quantization finding above), we chose not to
ship an approximated version under time pressure -- a scoping decision,
not a result.

We checked ARR's guidance before deciding whether a separate Ethical
considerations section was warranted: this work uses two long-established
public benchmarks (HotpotQA, GSM8K), two openly-released model
checkpoints, no human subjects, no crowdsourcing, and no new data
collection -- we are not aware of a concern specific to this paper beyond
what this section already covers.

\bibliography{references}
\bibliographystyle{acl_natbib}

\end{document}
